\documentclass[10pt,a4paper]{amsart}
\usepackage[margin=19mm]{geometry}
\usepackage{amsmath,amssymb,mathtools}
\usepackage[scr=rsfs,cal=euler]{mathalfa}
\usepackage{xfrac}
\usepackage[unicode,hidelinks,bookmarks=false]{hyperref}
\newcommand{\ie}{i.e.\ }
\newcommand{\cf}{cf.\ }
\newcommand{\resp}{respectively\ }
\newcommand{\Subsection}{\subsection}
\providecommand{\nobreakdash}{\nobreak\hbox{-}}
\setlength{\emergencystretch}{2em}
\multlinegap0pt
\usepackage{bm}
\usepackage{bbm}
\usepackage{dsfont}
\usepackage{xspace}
\usepackage{cancel}
\usepackage{mathtools}
\usepackage{amsthm}
\makeatletter
\@ifundefined{obs*}{\newtheorem*{obs*}{Remark}}{}
\makeatother
\title[An Extension of Major-Minor Mean Field Game Theory]{An Extension of Major-Minor Mean Field Game Theory}
\author{Agust\'in Mu\~noz Gonz\'alez}
\date{Source: arXiv:2603.16534v1 (CC BY 4.0)}
\begin{document}
\maketitle

\noindent\emph{Source and licence.} An Extension of Major-Minor Mean Field Game Theory. Version arXiv:2603.16534v1 (CC BY 4.0), distributed under the Creative Commons Attribution 4.0 International licence, as recorded in the arXiv licence metadata for this exact version and shown on the abstract page https://arxiv.org/abs/2603.16534v1. Full rights notice: licenses/arxiv-2603.16534-CC-BY-4.0.txt.\par\smallskip

\noindent\textit{Editorial scope.} Counted unit: the Obs* whose source label is \texttt{None} in the article (source0.tex lines 385--389, source label \texttt{None}), delivered together with the article's own complete original proof of it (source0.tex lines 391--457, one \texttt{proof} environment of 67 lines). No mathematics is written, repaired or re-derived: statement and proof are the article's own text line for line. Required context retained uncounted: Fokker-Planck (lines 299--308). Every definition, display and label of those lines is retained unchanged. Omitted from this package and not redistributed: 1--298, 309--384, 390--390, 458--763. Nothing inside a retained excerpt is omitted or altered: every retained body is delivered exactly as the article's lines are written, so no omission receipt of any kind accompanies this package. Citations: the retained excerpts carry no citation command at all, so no bibliography entry of the article is transcribed and no reference is redistributed. Reference adaptation: none was needed and none was made; every reference inside the delivered unit resolves inside it, so no unresolved reference or citation is produced. Printed slips kept literal: no printed word, symbol, hypothesis or number of the counted statement or of its proof was altered. Layout and class: the article's own class is \texttt{article} with packages including geometry, graphicx, amsmath, amsfonts, amsthm, authblk, blindtext, amssymb, tensor, epigraph, xspace, listings, stmaryrd, float; the wrapper is the lane's pinned \texttt{amsart} editorial wrapper at 10pt on A4, which restates the theorem-like environments the retained text needs and adds the article's own macro definitions that the retained text needs (none), with extra wrapper packages (bm, bbm, dsfont, xspace, cancel, mathtools). The delivered numbering is the unit's own. Assets: no retained span contains a figure environment, a graphics-inclusion command, a PDF-page inclusion, a table or a TikZ picture; no image file is copied into this package, the package declares no asset dependency and no image file is redistributed with it. Licence: version arXiv:2603.16534v1 of this article is distributed under the Creative Commons Attribution 4.0 International licence, as recorded in the arXiv licence metadata for this exact version and shown on the abstract page https://arxiv.org/abs/2603.16534v1. A full rights notice is delivered at licenses/arxiv-2603.16534-CC-BY-4.0.txt. Characters: every retained excerpt is pure ASCII, so the delivered engine, XeLaTeX with its default Latin Modern text font, reports no missing character.
\begin{equation}
    \label{Fokker-Planck}
    \begin{cases}
        \begin{aligned}
            \partial_t p_{u_1}(x,t) &= -A_1^* p_{u_1}(x,t)
            - \mathrm{div}\!\left(g_1\!\left(x,x_0(t),u_0(t), \Pi_t,u_1(x,t)\right)p_{u_1}(x,t)\right),\\[0.2cm]
            p_{u_1}(x,0) &= \omega(x),
        \end{aligned}
    \end{cases}
\end{equation}

\begin{obs*}[Martingale process $K_\Psi$]
The process $K_\Psi(x,t)$ appearing in the SHJB equation is the coefficient of the martingale part in the Doob--Meyer decomposition of $\Psi$ with respect to the filtration $\mathcal{F}^0$ generated by the Brownian motion of the dominating player $W_0$. More precisely, the value function $\Psi(x,t)$ satisfies a BSDE of the form
$$-d\Psi(x,t) = F(x,t,D\Psi)\,dt - K_\Psi(x,t)\,dW_0(t),$$
where $K_\Psi$ is the square-integrable integrand that ensures the adaptedness of the solution. The existence of $K_\Psi$ is guaranteed by the martingale representation theorem in the Brownian filtration $\mathcal{F}^0$.
\end{obs*}

\begin{proof}
We apply the stochastic maximum principle. For any perturbation \(\tilde{u}_1 \in \mathcal{A}_1\),
\begin{equation}
\label{Condicion optimalidad}
    \begin{aligned}
    0 &= \left.\frac{d}{d\theta}\right|_{\theta=0} J_1(\hat{u}_1+\theta\tilde{u}_1,x_0,u_0,\Pi).
    \end{aligned}
\end{equation}
Rewriting the expectation in terms of the conditional density $p_{\hat{u}_1}(x,t)$ of the state $x_1$ and differentiating, we obtain
\begin{align*}
0 &= \mathbb{E}\Bigg[\int_0^T\int_{\mathbb{R}^{n_1}} \tilde{p}(x,t)\, f_1(x,x_0(t),u_0(t),\Pi_t,\hat{u}_1(x,t)) \, dx\,dt \\
&\quad + \int_0^T\int_{\mathbb{R}^{n_1}} p_{\hat{u}_1}(x,t)\, f_{1,u_1}(x,x_0(t),u_0(t),\Pi_t,\hat{u}_1(x,t)) \, \tilde{u}_1(x,t)\, dx\,dt \\
&\quad + \int_{\mathbb{R}^{n_1}} \tilde{p}(x,T)\, h_1(x,x_0(T),u_0(T),\Pi_T)\, dx \Bigg],
\end{align*}
where \(\tilde{p}=\left.\tfrac{d}{d\theta}\right|_{\theta=0} p_{\hat{u}_1+\theta\tilde{u}_1}\).

Differentiating with respect to $\theta$ in the Fokker--Planck equation \eqref{Fokker-Planck}, $\tilde{p}$ satisfies
    \begin{equation*}
        \begin{cases}
            \begin{aligned}
                \frac{\partial \tilde{p}}{\partial t} &= -A_1^*\tilde{p}(x,t)-\mathrm{div}\left(\tilde{u}_1(x,t)g_{1,u_1}\left(x,x_0(t),u_0(t),\Pi_t,\hat{u}_1(x,t)\right)p_{\hat{u}_1}(x,t)\right)\\
                &\quad - \mathrm{div}\left(g_1\left(x,x_0(t),u_0(t),\Pi_t,\hat{u}_1(x,t)\right)\tilde{p}(x,t)\right),\\
                \tilde{p}(x,0) &= 0.
            \end{aligned}
        \end{cases}
    \end{equation*}

We introduce the adjoint process $\Psi$ as the solution of the BSDE
    \[
        \left\{
        \begin{aligned}
            -\partial_t \Psi &= \Big( f_1(x, x_0(t), u_0(t), \Pi_t, u_1) +  D\Psi(x,t)\, g_1(x,x_0(t),u_0(t),\Pi_t,u_1) - A_1 \Psi(x,t) \Big)\, dt \\
            &\quad - K_\Psi(x,t)\, dW_0(t), \\
            \Psi(x,T) &= h_1 \left( x, x_0(T), u_0(T), \Pi_T \right).
        \end{aligned}
        \right.
    \]

Consider the inner product
    \begin{align*}
        & d\int_{\mathbb{R}^{n_1}}\tilde{p}(x,t)\Psi(x,t)dx \\
        &= \int_{\mathbb{R}^{n_1}} \left\{-A_1^*\tilde{p}(x,t)-\mathrm{div}\left[\tilde{u}_1(x,t)g_{1,u_1}(x,x_0(t),u_0(t),\Pi_t,\hat{u}_1(x,t))p_{\hat{u}_1}(x,t)\right]\right.\\
        &\hspace{2cm} \left. -\mathrm{div}\left[g_1(x,x_0(t),u_0(t),\Pi_t,\hat{u}_1(x,t))\tilde{p}(x,t)\right] \right\}\Psi(x,t)dxdt\\
        &\hspace{1cm} -\int_{\mathbb{R}^{n_1}} \tilde{p}(x,t)\left\{ f_1(x, x_0(t), u_0(t), \Pi_t, u_1(x,t)) \right.\\
        &\hspace{1cm}+\left. D\Psi(x,t)g_1(x,x_0(t),u_0(t),\Pi_t,u_1(x,t)) - A_1 \Psi(x,t) \right\} dxdt \\
        &\hspace{1cm} + \int_{\mathbb{R}^{n_1}} \tilde{p}(x,t)K_\Psi(x,t) dxdW_0(t)\\
        &=\int_{\mathbb{R}^{n_1}} \left( \tilde{u}_1(x,t)g_{1,u_1}(x,x_0(t),u_0(t),\Pi_t,\hat{u}_1(x,t))p_{\hat{u}_1}(x,t) \right)D\Psi(x,t)dxdt\\
        & -\int_{\mathbb{R}^{n_1}}\tilde{p}(x,t) f_1(x, x_0(t), u_0(t), \Pi_t, u_1(x,t)) dxdt+ \int_{\mathbb{R}^{n_1}} \tilde{p}(x,t)K_\Psi(x,t) dxdW_0(t).
    \end{align*}

Integrating over $[0,T]$ and taking expectations on both sides we get
    \begin{align*}
        & \mathbb{E}\left[\int_{\mathbb{R}^{n_1}}\tilde{p}(x,T)h_1(x,x_0(T),u_0(T),\Pi_T)dx\right] \\
        &=\mathbb{E}\left[\int_0^T\int_{\mathbb{R}^{n_1}}\left( \tilde{u}_1(x,t)g_{1,u_1}(x,x_0(t),u_0(t),\Pi_t,\hat{u}_1(x,t))p_{\hat{u}_1}(x,t) \right)D\Psi(x,t)dxdt \right.\\
        &\hspace{1cm} \left.-\int_0^T\int_{\mathbb{R}^{n_1}}\tilde{p}(x,t) f_1(x, x_0(t), u_0(t), \Pi_t, u_1(x,t)) dxdt\right],
    \end{align*}
where the term involving $K_\Psi$ vanishes upon taking expectations since it is a martingale term with zero expectation. Combining this with equation \eqref{Condicion optimalidad}, we obtain
    \begin{align*}
        0 &= \mathbb{E}\left[\int_0^T\int_{\mathbb{R}^{n_1}} \tilde{u}_1(x,t)\left[g_{1,u_1}(x,x_0(t),u_0(t),\Pi_t,\hat{u}_1(x,t))D\Psi(x,t) \right. \right. \\
        &+ \left. \left. f_{1,u_1}(x,x_0(t),u_0(t),\Pi_t,\hat{u}_1(x,t))\right] p_{\hat{u}_1}(x,t)dxdt\right].
    \end{align*}
Recall that $p_{\hat{u}_1}(\cdot,t)$ is a conditional probability density function and hence nonnegative, and $\tilde{u}_1$ is an arbitrary Markovian control. Therefore $\hat{u}_1$ is optimal only if
    $$g_{1,u_1}(x,x_0(t),u_0(t),\Pi_t,\hat{u}_1(x,t))D\Psi(x,t)+ f_{1,u_1}(x,x_0(t),u_0(t),\Pi_t,\hat{u}_1(x,t))=0,\quad a.e.\,(x,t).$$
This is equivalent to
    $$L_{u_1}(x,x_0(t),u_0(t),\Pi_t,\hat{u}_1(x,t), D\Psi(x,t))=0,\quad a.e.\,(x,t),$$
which provides a necessary condition for the minimization problem. Since the minimizer is assumed to be attained at $\hat{u}_1$, which depends on $x, x_0, u_0, \Pi$ and $D\Psi$, we obtain the SHJB equation.
\end{proof}

\end{document}
